%% =====================================================================
%%  Unit 10 — How much is it?  (Faire des achats)
%%  Pistes : 10.01–10.04 dialogues · 10.05–10.06 vocabulaire · 10.07–10.08 grammaire
%%           10.09 prononciation · 10.10–10.13 exercices (textes en annexe)
%% =====================================================================

\newcommand{\otIllusUnX}{\otIllusDuo{\$15/104/86/15/0.10,sale/190/92/20/0.09,shoes/122/14/13/0.08,
  \$50/150/8/11/0.08,coat/196/22/12/0.07,cash/100/50/12/0.07}{lea}{How much is it?}{emma}{It's fifteen dollars!}}

\unite[How much\\is it?]{10}{How much is it?}{Faire des achats}{otIllusUnX}

\situation{Épisode 10}{Fin octobre : il commence à faire froid à New York, et Léa n'a pas de
vêtements chauds. Samedi, elle fait les boutiques. Le soir, Sam l'accompagne au supermarché : il
manque presque tout dans son réfrigérateur !}

\begin{objectifs}
\item demander un prix et le comprendre : \en{How much is it?} ;
\item acheter un vêtement, l'essayer et payer ;
\item employer \en{how much} et \en{how many}, \en{some} et \en{any} ;
\item distinguer les noms dénombrables et indénombrables ;
\item ne plus confondre \en{fifteen} et \en{fifty}.
\end{objectifs}

\auprogramme
  {Demander un prix, acheter, payer}
  {\en{How much}, \en{how many}\par\vspace{0.6mm}\en{some}, \en{any}}
  {Vêtements, magasins, prix : 29 expressions}
  {\en{-teen} ou \en{-ty} ? \en{thirteen}, \en{thirty}}
  {Dollars, cents et taxes}

\semaine{casque/otNavy/Dialogues/Écouter et lire les deux scènes,
  bulle/otGreenDark/Vocabulaire/Apprendre les 29 expressions,
  blocs/otViolet/Grammaire/how much et some any,
  sons/otOrangeDark/Sons et culture/teen ou ty et les dollars,
  crayon/otNavy/Exercices/12 exercices corrigés,
  coche/otGreenDark/Bilan/L'essentiel et le test express}

\clearpage
%% ---------------------------------------------------------------------
\rubrique{ecouter}{Les courses}

\begin{dialogue}{1}{How much is this sweater?}{Samedi, 11 heures, dans un magasin de vêtements de
  Brooklyn. Une vendeuse s'approche de Léa.}{10.01}{10.02}
\dit{vendeuse}{Hi! \hl{Can I help you}?}{Bonjour ! Je peux vous aider ?}
\dit{lea}{Yes, please. \hl{How much is} this sweater?\note{1}}{Oui, s'il vous plaît. Combien coûte ce pull ?}
\dit{vendeuse}{It's forty dollars.}{Il coûte quarante dollars.}
\dit{lea}{And \hl{how much are} these gloves?}{Et combien coûtent ces gants ?}
\dit{vendeuse}{They're fifteen dollars.}{Ils coûtent quinze dollars.}
\dit{lea}{Fifteen or fifty?\note{2}}{Quinze ou cinquante ?}
\dit{vendeuse}{Fifteen: one, five.}{Quinze : un, cinq.}
\dit{lea}{Great! \hl{Can I try on} the sweater?}{Super ! Je peux essayer le pull ?}
\dit{vendeuse}{Of course. The fitting room is over there.}{Bien sûr. La cabine d'essayage est là-bas.}
\didascalie{Cinq minutes plus tard.}
\dit{lea}{\hl{I'll take it}! And the gloves too.\note{3}}{Je le prends ! Et les gants aussi.}
\dit{vendeuse}{OK. That's fifty-five dollars, \hl{plus tax}.}{D'accord. Ça fait cinquante-cinq dollars, plus les taxes.}
\dit{lea}{\hl{Can I pay by card}?}{Je peux payer par carte ?}
\dit{vendeuse}{Sure!}{Bien sûr !}
\end{dialogue}

\begin{comprendre}
\pointnote{1} \en{\textbf{How much is this sweater?}} « Combien coûte ce pull ? ». Pour une seule
  chose : \en{How much is…?} ; pour plusieurs : \en{How much are…?} On répond \en{It's…} ou
  \en{They're…} + le prix.
\pointnote{2} \en{\textbf{Fifteen or fifty?}} Une confusion classique : \en{fifteen} (15) est accentué
  sur la fin, \en{fifty} (50) sur le début. Voir la rubrique Prononciation.
\pointnote{3} \en{\textbf{I'll take it!}} « Je le prends ! », la formule pour acheter. \en{Can I…?}
  (« Est-ce que je peux… ? ») sert à demander la permission : \en{Can I try it on?},
  \en{Can I pay by card?} Vous étudierez \en{can} aux unités 11 et 12.
\end{comprendre}

\begin{dialogue}{2}{We need some milk}{Samedi, 18 heures, au supermarché. Léa et Sam font les
  courses ensemble.}{10.03}{10.04}
\dit{sam}{OK, what do we need?}{Bon, qu'est-ce qu'il nous faut ?}
\dit{lea}{We need \hl{some} milk, \hl{some} bread and \hl{some} eggs.\note{4}}{Il nous faut du lait, du pain et des œufs.}
\dit{sam}{\hl{How many} eggs?}{Combien d'œufs ?}
\dit{lea}{Six. And some cheese.}{Six. Et du fromage.}
\dit{sam}{Sorry, there \hl{isn't any} French cheese here!\note{5}}{Désolé, il n'y a pas de fromage français ici !}
\dit{lea}{Oh no! Are there \hl{any} apples?}{Oh non ! Il y a des pommes ?}
\dit{sam}{Yes, there are some apples here. They're cheap!}{Oui, il y a des pommes ici. Elles ne sont pas chères !}
\dit{lea}{Great. And coffee, of course!}{Super. Et du café, bien sûr !}
\dit{sam}{\hl{How much} coffee do you need?\note{6}}{Il te faut combien de café ?}
\dit{lea}{Two bags. I drink \hl{a lot of} coffee!}{Deux paquets. Je bois beaucoup de café !}
\end{dialogue}

\begin{comprendre}
\pointnote{4} \en{\textbf{some milk}, \textbf{some eggs}} \en{some} traduit « du », « de la »,
  « des » dans une phrase affirmative : \en{some milk} (du lait), \en{some eggs} (des œufs).
\pointnote{5} \en{\textbf{there isn't any cheese}} Dans une phrase négative ou une question, \en{some}
  devient \en{any} : « il n'y a pas de fromage », « il y a des pommes ? ».
\pointnote{6} \en{\textbf{How much coffee?}} Pour une quantité, \en{how much} s'emploie avec un nom
  qu'on ne peut pas compter (\en{coffee}, \en{milk}) et \en{how many} avec un nom pluriel
  (\en{eggs}, \en{apples}). \en{a lot of} = « beaucoup de ».
\end{comprendre}

%% ---------------------------------------------------------------------
\rubrique{mots}{Les mots de l'unité}

Voici les 29 expressions de l'unité. Attention : \en{pants}, \en{jeans}, \en{shoes} et \en{gloves}
sont toujours au pluriel en anglais, même pour un seul pantalon : \en{These pants are nice.}

\begin{vocabulaire}{Les vêtements}{10.05}
\categorie{Le haut et le bas}
\mot{clothes}{kloʊðz}{les vêtements}
\mot{a sweater}{ə ˈswetɚ}{un pull}
\mot{a T-shirt}{ə ˈtiːʃɝːt}{un tee-shirt}
\mot{a shirt}{ə ˈʃɝːt}{une chemise}
\mot{a dress}{ə ˈdres}{une robe}
\mot{a skirt}{ə ˈskɝːt}{une jupe}
\mot{pants}{pænts}{un pantalon \registre{américain}}
\mot{jeans}{dʒiːnz}{un jean}
\categorie{Pour sortir}
\mot{a jacket}{ə ˈdʒækɪt}{une veste, un blouson}
\mot{a coat}{ə ˈkoʊt}{un manteau}
\mot{shoes}{ʃuːz}{des chaussures}
\mot{gloves}{ɡlʌvz}{des gants}
\mot{a scarf}{ə ˈskɑːrf}{une écharpe}
\mot{a hat}{ə ˈhæt}{un chapeau, un bonnet}
\end{vocabulaire}

\begin{vocabulaire}{Acheter et payer}{10.06}
\categorie{Dans le magasin}
\mot{a store}{ə ˈstɔːr}{un magasin}
\mot{Can I help you?}{kən aɪ ˈhelp juː}{Je peux vous aider ?}
\mot{What size?}{ˌwʌt ˈsaɪz}{Quelle taille ?}
\mot{Can I try it on?}{kən aɪ ˌtraɪ ɪt ˈɑːn}{Je peux l'essayer ?}
\mot{I'll take it.}{aɪl ˈteɪk ɪt}{Je le prends.}
\mot{on sale}{ɑːn ˈseɪl}{en promotion, en solde}
\categorie{Les prix}
\mot{How much is it?}{ˌhaʊ ˈmʌtʃ ɪz ɪt}{Combien ça coûte ?}
\mot{How much are they?}{ˌhaʊ ˈmʌtʃ ɑːr ðeɪ}{Combien coûtent-ils ?}
\mot{It's twenty dollars.}{ɪts ˌtwenti ˈdɑːlɚz}{Ça coûte vingt dollars.}
\mot{a cent}{ə ˈsent}{un cent (centième de dollar)}
\mot{cheap · expensive}{tʃiːp · ɪkˈspensɪv}{bon marché · cher}
\mot{Can I pay by card?}{kən aɪ ˌpeɪ baɪ ˈkɑːrd}{Je peux payer par carte ?}
\mot{cash}{kæʃ}{en espèces, du liquide}
\mot{money}{ˈmʌni}{l'argent}
\mot{a lot of}{ə ˈlɑːt əv}{beaucoup de}
\end{vocabulaire}

%% ---------------------------------------------------------------------
\rubrique{grammaire}{How much, how many, some et any}

\otIntertitre[otViolet]{1 · Dénombrable ou indénombrable ?}
\begin{paire}
\begin{regle}[Dénombrables : on peut les compter]
\exen{an apple, two apples}{une pomme, deux pommes}
\exen{an egg, six eggs}{un œuf, six œufs}
\exen{a dollar, ten dollars}{un dollar, dix dollars}
Ils ont un singulier (\en{a}, \en{an}) et un pluriel.
\end{regle}
\begin{regle}[Indénombrables : on ne les compte pas]
\exen{milk, water, coffee}{le lait, l'eau, le café}
\exen{bread, cheese, rice}{le pain, le fromage, le riz}
\exen{money, information}{l'argent, les informations}
Ni \en{a}, ni pluriel : \en{some bread}, \en{a lot of money}.
\end{regle}
\end{paire}

\otIntertitre[otViolet]{2 · \en{How much} ou \en{how many} ?}
\begin{center}\otTableau
\begin{tabular}{@{}l@{\hspace{8mm}}>{\enOn}l<{\enOff}@{\hspace{8mm}}l@{}}
\otTH{Pour demander} & \otTH{Exemple} & \otTH{Français}\\\midrule
un prix                          & How much is it? / How much are they? & Combien ça coûte ?\\
une quantité (indénombrable)     & How much coffee do you need?         & Combien de café ?\\
un nombre (dénombrable pluriel)  & How many eggs do you need?           & Combien d'œufs ?\\\bottomrule
\end{tabular}\end{center}

\blocaudio{10.07}{Écoutez et répétez}{Six questions de prix ou de quantité, et leur réponse.}

\otIntertitre[otViolet]{3 · \en{some} ou \en{any} ?}
\begin{center}\otTableau
\begin{tabular}{@{}l@{\hspace{8mm}}>{\enOn}l<{\enOff}@{\hspace{8mm}}l@{}}
\otTH{Phrase} & \otTH{Exemple} & \otTH{Français}\\\midrule
affirmative & We need \textbf{some} milk. There are \textbf{some} apples. & du lait, des pommes\\
négative    & There isn't \textbf{any} cheese. I don't have \textbf{any} money. & pas de fromage, pas d'argent\\
question    & Are there \textbf{any} apples? Do you need \textbf{any} bread? & des pommes ? du pain ?\\\bottomrule
\end{tabular}\end{center}
{\fontsize{10}{13}\selectfont Exception utile : pour \textbf{proposer} ou \textbf{demander} quelque chose,
on garde \en{some} dans la question : \en{Would you like some coffee?} (unité 11).\par}

\blocaudio{10.08}{Écoutez et répétez}{Six phrases avec \en{some} ou \en{any}.}

\begin{attention}[Les pièges des francophones]
\fauxjuste{How much eggs?}{How many eggs?}
\fauxjuste{a bread}{some bread / a loaf of bread}
\fauxjuste{There isn't some cheese.}{There isn't any cheese.}
\fauxjuste{informations}{information}
\fauxjuste{How much costs it?}{How much is it?}
\end{attention}

%% ---------------------------------------------------------------------
\rubrique[-teen ou -ty]{sons}{\en{-teen} ou \en{-ty} ? \en{thirteen} ou \en{thirty} ?}

Au téléphone, au marché ou à la caisse, confondre 15 et 50 peut coûter cher ! Deux indices
permettent de les distinguer : l'\textbf{accent} (sur \en{-teen}, mais \textbf{avant} \en{-ty}) et
la \textbf{fin du mot} (le \phon{n} de \en{-teen} s'entend nettement).

\begin{paires}{10.09}{Écoutez chaque paire, puis répétez en exagérant l'accent : \en{thir\textbf{TEEN}},
\en{\textbf{THIR}ty}.}
\sonpaire{thir\acc{teen}}{θɝːˈtiːn}{13}{\acc{thir}ty}{ˈθɝːti}{30}
\sonpaire{four\acc{teen}}{ˌfɔːrˈtiːn}{14}{\acc{for}ty}{ˈfɔːrti}{40}
\sonpaire{fif\acc{teen}}{ˌfɪfˈtiːn}{15}{\acc{fif}ty}{ˈfɪfti}{50}
\sonpaire{six\acc{teen}}{ˌsɪkˈstiːn}{16}{\acc{six}ty}{ˈsɪksti}{60}
\sonpaire{seven\acc{teen}}{ˌsevənˈtiːn}{17}{\acc{sev}enty}{ˈsevənti}{70}
\sonpaire{eigh\acc{teen}}{ˌeɪˈtiːn}{18}{\acc{eigh}ty}{ˈeɪti}{80}
\sonpaire{nine\acc{teen}}{ˌnaɪnˈtiːn}{19}{\acc{nine}ty}{ˈnaɪnti}{90}
\sonpaire{fif\acc{teen} dollars}{ˌfɪfˈtiːn ˈdɑːlɚz}{\$15}{\acc{fif}ty dollars}{ˈfɪfti ˈdɑːlɚz}{\$50}
\end{paires}

\begin{astuce}[En cas de doute]
Demandez de préciser chiffre par chiffre, comme Léa : \en{Fifteen or fifty?} — \en{Fifteen: one,
five.} C'est parfaitement naturel, même entre anglophones.
\end{astuce}

%% ---------------------------------------------------------------------
\rubrique{culture}{Dollars, cents et taxes}

\begin{paire}
\begin{culture}[La monnaie]
Un dollar (\en{\$1}) vaut 100 \en{cents}. Les pièces ont un surnom : \en{a nickel} (5 cents),
\en{a dime} (10 cents) et \en{a quarter} (25 cents). Les billets, de 1 à 100 dollars, ont tous la
même taille : vérifiez bien le chiffre avant de payer !
\end{culture}
\begin{culture}[La taxe en plus]
Aux États-Unis, les prix affichés s'entendent \textbf{hors taxe}. La taxe sur les ventes
(\en{sales tax}) est ajoutée à la caisse : à New York, elle approche 9 \%. Un pull à 40 dollars
coûte donc un peu plus de 43 dollars.
\end{culture}
\end{paire}

\begin{culture}[Lire un prix]
\en{\$12.99} se lit \en{twelve ninety-nine} ou \en{twelve dollars and ninety-nine cents}. Le point
sert de séparateur décimal et la virgule sépare les milliers : \en{\$2,400} = deux mille quatre
cents dollars. Et attention : au Royaume-Uni, \en{pants} désigne… un slip ! Le pantalon s'y dit
\en{trousers}.
\end{culture}

%% ---------------------------------------------------------------------
\rubrique{pratique}{Exercices}

\exercice[ecrire,niveau=1]{Les vêtements}{Écrivez le mot anglais.}
\begin{questions}[2][french]
\item un pull → \enb{\trou[30mm]{a sweater}}
\item une robe → \enb{\trou[30mm]{a dress}}
\item un manteau → \enb{\trou[30mm]{a coat}}
\item des chaussures → \enb{\trou[30mm]{shoes}}
\item une écharpe → \enb{\trou[30mm]{a scarf}}
\item un pantalon (US) → \enb{\trou[30mm]{pants}}
\item une jupe → \enb{\trou[30mm]{a skirt}}
\item des gants → \enb{\trou[30mm]{gloves}}
\end{questions}

\exercice[ecrire,niveau=1]{\en{a}, \en{an} ou \en{some} ?}{Complétez : \en{a} ou \en{an} devant un
dénombrable singulier, \en{some} devant un indénombrable ou un pluriel.}
\begin{questions}[2]
\item \trou[16mm]{an} apple
\item \trou[16mm]{some} milk
\item \trou[16mm]{some} bread
\item \trou[16mm]{an} egg
\item \trou[16mm]{some} eggs
\item \trou[16mm]{some} coffee
\item \trou[16mm]{a} sweater
\item \trou[16mm]{some} money
\end{questions}

\exercice[ecrire,niveau=2]{\en{How much} ou \en{how many} ?}{Complétez.}
\begin{questions}
\item \trou[30mm]{How much} is this coat? — It's ninety dollars.
\item \trou[30mm]{How many} eggs do you need? — Six.
\item \trou[30mm]{How much} milk do we need? — One bottle.
\item \trou[30mm]{How much} are these shoes? — They're sixty dollars.
\item \trou[30mm]{How many} apples are there? — There are ten.
\item \trou[30mm]{How much} money do you have? — Twenty dollars.
\end{questions}

\exercice[ecrire,niveau=2]{\en{some} ou \en{any} ?}{Complétez.}
\begin{questions}
\item We need \trou[16mm]{some} bread.
\item There isn't \trou[16mm]{any} milk in the fridge.
\item Are there \trou[16mm]{any} apples?
\item I'd like \trou[16mm]{some} coffee, please.
\item Sorry, we don't have \trou[16mm]{any} French cheese.
\item There are \trou[16mm]{some} eggs on the table.
\end{questions}

\exercice[lire,ecrire,niveau=1]{Les prix}{Écrivez chaque prix en lettres, comme on le dit.}
\begin{questions}[2]
\item \$12 → \trou[44mm]{twelve dollars}
\item \$30 → \trou[44mm]{thirty dollars}
\item \$99 → \trou[44mm]{ninety-nine dollars}
\item 75¢ → \trou[44mm]{seventy-five cents}
\item \$15.50 → \trou[44mm]{fifteen fifty}
\item \$2.99 → \trou[44mm]{two ninety-nine}
\end{questions}

\exercice[lire,niveau=1]{Au magasin}{Remettez le dialogue dans l'ordre : écrivez les lettres.}
\begin{anglais}
\begin{questions}[1][english]
\item[a.] It's thirty-five dollars.
\item[b.] Hi! Can I help you?
\item[c.] Yes. How much is this scarf?
\item[d.] Sure. That's thirty-five dollars, plus tax.
\item[e.] Great! I'll take it. Can I pay by card?
\end{questions}
\end{anglais}
\par\vspace{2mm}\noindent{\otSemi\fontsize{10}{12}\selectfont Ordre :}\hspace{4mm}%
\trou[9mm]{b}\hspace{3mm}\trou[9mm]{c}\hspace{3mm}\trou[9mm]{a}\hspace{3mm}\trou[9mm]{e}\hspace{3mm}\trou[9mm]{d}

\exercice[ecouter,niveau=1,piste=10.10]{15 ou 50 ?}{Écoutez. Pour chaque numéro, cochez le nombre
entendu.}
\begin{questions}[4]
\item \choix{13}\quad\choix{30}
\item \choix{14}\quad\choix{40}
\item \choix{15}\quad\choix{50}
\item \choix{16}\quad\choix{60}
\item \choix{17}\quad\choix{70}
\item \choix{18}\quad\choix{80}
\item \choix{19}\quad\choix{90}
\item \choix{\$15}\quad\choix{\$50}
\end{questions}

\exercice[ecouter,ecrire,niveau=2,piste=10.11]{Combien ça coûte ?}{Emma fait les magasins. Écoutez
et notez le prix de chaque article.}
\begin{center}\otTableau
\begin{tabular}{@{}l@{\hspace{14mm}}l@{}}
\otTH{Article} & \otTH{Prix}\\\midrule
{\otSemi 1 · \enb{a red dress}}        & \trou[26mm]{\$45}\\[1mm]
{\otSemi 2 · \enb{black shoes}}        & \trou[26mm]{\$80}\\[1mm]
{\otSemi 3 · \enb{a T-shirt}}          & \trou[26mm]{\$13}\\[1mm]
{\otSemi 4 · \enb{a winter coat}}      & \trou[26mm]{\$99}\\\bottomrule
\end{tabular}\end{center}

\exercice[ecouter,ecrire,niveau=2,piste=10.12]{Dictée à trous}{Sam et Léa préparent la liste de
courses. Écoutez et complétez.}
\begin{anglais}
\noindent\begin{tabular}{@{}l@{\hspace{6mm}}l@{}}
\otNomDit{sam} & Is there \trou{any} milk in the fridge?\\[1.8mm]
\otNomDit{lea} & No, there \trou{isn't}. We need \trou{some} milk.\\[1.8mm]
\otNomDit{sam} & How \trou{many} eggs do we need?\\[1.8mm]
\otNomDit{lea} & Twelve. And \trou{some} bread, please.\\[1.8mm]
\otNomDit{sam} & How \trou{much} money do we have?\\[1.8mm]
\otNomDit{lea} & Forty \trou{dollars}. That's OK!\\
\end{tabular}
\end{anglais}

\exercice[ecrire,niveau=2]{Traduisez}{Traduisez en anglais.}
\begin{questions}[1][french]
\item Combien coûte ce manteau ? — Il coûte cent dollars.\lignes{1}
\item Je peux l'essayer ?\lignes{1}
\item Il n'y a pas de lait.\lignes{1}
\item Combien d'œufs y a-t-il ?\lignes{1}
\item Il nous faut du pain et du fromage.\lignes{1}
\item Ces chaussures sont chères.\lignes{1}
\item Je le prends. Je peux payer en espèces ?\lignes{1}
\item Je bois beaucoup de café.\lignes{1}
\end{questions}

\exercice[ecrire,niveau=3]{Au magasin}{Écrivez un petit dialogue (six répliques) entre un vendeur et
vous : demandez un prix, essayez, achetez, payez.}
\lignes{6}

\exercice[ecouter,parler,niveau=3,piste=10.13]{À vous de parler !}{Vous êtes dans un magasin de
vêtements. Répondez à la vendeuse \textbf{à voix haute} pendant chaque pause.}
\begin{center}\otTableau
\begin{tabular}{@{}>{\otSemi\fontsize{8.6}{10}\selectfont\color{otNavy}}l@{\hspace{3mm}}>{\enOn}p{70mm}<{\enOff}@{\hspace{6mm}}p{78mm}@{}}
 & \otTH{La vendeuse dit\ldots} & \otTH{Vous\ldots}\\\midrule
1. & Hi! Can I help you? & dites que vous cherchez un pull, et demandez le prix de celui-ci.\\
2. & It's sixty-five dollars. & demandez si vous pouvez l'essayer.\\
3. & Of course! What size? & donnez votre taille (\en{small}, \en{medium} ou \en{large}).\\
4. & Here you are. Do you like it? & dites que vous le prenez.\\
5. & Great! Cash or card? & dites comment vous payez.\\\bottomrule
\end{tabular}\end{center}

%% ---------------------------------------------------------------------
\rubrique{bilan}{Faites le point}

\begin{essentiel}
\begin{minipage}[t]{0.55\linewidth}
\otEssTitre{Les phrases clés}
\renewcommand{\arraystretch}{1.32}
\begin{tabular}{@{}>{\enOn\otSemi\fontsize{9.4}{11}\selectfont}l<{\enOff}@{\hspace{3mm}}>{\fontsize{9.2}{11}\selectfont\color{white!78!otNavy}}l@{}}
How much is it?             & Combien ça coûte ?\\
It's fifteen dollars.       & Ça coûte 15 dollars.\\
Can I try it on?            & Je peux l'essayer ?\\
I'll take it.               & Je le prends.\\
Can I pay by card?          & Je peux payer par carte ?\\
We need some milk.          & Il nous faut du lait.\\
\end{tabular}
\end{minipage}\hfill
\begin{minipage}[t]{0.4\linewidth}
\otEssTitre{Quantités}
{\enb{\fontsize{9.4}{12}\selectfont how many + pluriel · how much + indénombrable}\par}
\otEssTitre{some / any}
{\fontsize{9.2}{11.6}\selectfont \en{some} : affirmation · \en{any} : négation, question\par}
\otEssTitre{Prononciation}
{\fontsize{9.2}{11.6}\selectfont \en{fif\textbf{teen}} (15) · \en{\textbf{fif}ty} (50)\par}
\end{minipage}
\end{essentiel}

\begin{jesais}
\item demander un prix et comprendre la réponse ;
\item essayer, acheter et payer un vêtement dans un magasin ;
\item nommer une quinzaine de vêtements ;
\item choisir entre \en{how much} et \en{how many} ;
\item employer \en{some} et \en{any}, et reconnaître un indénombrable ;
\item distinguer \en{fifteen} et \en{fifty}, à l'écoute comme à l'oral.
\end{jesais}

\clearpage
\begin{testexpress}
\qcm[b]{« Combien ça coûte ? » se dit :}{\en{How much costs it?}}{\en{How much is it?}}{\en{How many is it?}}
\qcm[c]{Complétez : \en{How \rule[-0.6mm]{9mm}{0.5pt} eggs do you need?}}{\en{much}}{\en{lot}}{\en{many}}
\qcm[a]{Complétez : \en{How \rule[-0.6mm]{9mm}{0.5pt} milk do we need?}}{\en{much}}{\en{many}}{\en{some}}
\qcm[b]{Complétez : \en{There isn't \rule[-0.6mm]{9mm}{0.5pt} cheese.}}{\en{some}}{\en{any}}{\en{a}}
\qcm[a]{« Il nous faut du pain » se dit :}{\en{We need some bread.}}{\en{We need a bread.}}{\en{We need any bread.}}
\qcm[c]{« Je le prends » (au magasin) se dit :}{\en{I take him.}}{\en{I'm taking it.}}{\en{I'll take it.}}
\qcm[b]{50 se dit :}{\en{fifteen}}{\en{fifty}}{\en{five-ten}}
\qcm[a]{Dans \en{fifteen}, l'accent tombe sur :}{\en{teen}}{\en{fif}}{les deux syllabes}
\qcm[c]{« un pantalon » se dit aux États-Unis :}{\en{a pant}}{\en{trousers}}{\en{pants}}
\qcm[b]{Aux États-Unis, la taxe sur les ventes est :}{comprise dans le prix affiché}{ajoutée à la caisse}{remboursée à la sortie}
\end{testexpress}

\noindent{\fontsize{9.6}{12}\selectfont\color{otMuted}Niveau visé (CECRL, niveau A1) : « peut se
débrouiller avec les nombres, les quantités, l'heure et les prix ; peut faire des achats simples ».
Les corrigés de l'unité sont page~\pageref{cor-u10-1} ; les textes des pistes d'exercices,
page~\pageref{trans-10.10}.\par}
